\vspace{-2mm}\section{Related Work}\vspace{-1mm}

As LLMs shift from instruction-following to long-horizon deliberation, self-verification, and test-time computation, existing surveys have focused on pretraining, alignment, and factual reliability limits~\citep{QIN2025101118,zhao2023survey,huang2023hallucination,ji2023hallucination}. Early prompting strategies like Chain-of-Thought and its variants (Tree/Graph-of-Thought) demonstrated that intermediate reasoning improves multi-step problem solving~\citep{wei2022chain,kojima2022large,wang2022self,yao2023tree,besta2024graph}. Recent works explore inference-time scaling, process supervision, and reinforcement-learning-driven reasoning~\citep{chen2025towards,li2025system1system2,deepseek2025r1,snell2024scaling,muennighoff2025s1}, alongside self-improvement and reflection feedback loops~\citep{madaan2023selfrefine,shinn2023reflexion}. However, while most literature treats reasoning as an internal, text-centered capability, we focus on how it becomes fully operational when integrated with tools, persistent states, and verifiable workspace changes~\citep{jeong2024survey,chang2024survey,asante2026evaluation,masterman2024landscape,chen2025ai4researchsurveyartificialintelligence}.

While \citet{wang2024survey} and \citet{xi2023rise} overview broad agent architectures, specialized surveys highlight the planning, memory, and reliability safeguards required for long-horizon tasks~\citep{wei2025plangenllms,zhang2025survey,lin2025llm}. Cognitive frameworks analyze how memory and planning sustain this long-term behavior~\citep{sumers2023cognitive,park2023generative}. To execute actions, \citet{yao2022react} integrate reasoning with external steps, while \citet{schick2023toolformer}, \citet{qin2023toolllm}, \citet{shen2023hugginggpt}, and \citet{patil2023gorilla} investigate API invocation, task routing, and tool utilization. Additionally, \citet{wang2023voyager} demonstrates open-ended skill accumulation. Together, these foundational works establish the traditional "observe-plan-act" loop. In contrast, we highlight a transition: agents are evolving from simple tool callers into operators of durable workspaces containing files, sessions, and permissions~\citep{chowa2026language,chen2025tool,zhai2025survey}.

Further, interactive benchmarks are essential for evaluating whether agents can complete tasks in realistic environments. \citet{liu2023agentbench} and \citet{mialon2023gaia} evaluate general agent abilities and assistant-style problem solving. Earlier browser-assisted systems like WebGPT connected language models with web navigation and human feedback~\citep{nakano2022webgpt}. Web-based environments such as WebShop, Mind2Web, and WebArena then test agents through interactive shopping, website navigation, and realistic web tasks~\citep{yao2022webshop,deng2023mind2web,zhou2023webarena}. OSWorld and WorkArena extend this direction to desktop operating systems and enterprise workflows~\citep{xie2024osworld,drouin2024workarena}. In software engineering, SWE-bench, SWE-agent, and OpenHands emphasize repositories, tests, execution feedback, and agent-computer interfaces for development tasks~\citep{jimenez2023swebench,yang2024sweagent,wang2024openhands}. Tool-learning benchmarks and datasets such as APIBank, ToolAlpaca, ToolLLM, ToolSandbox, and 
-bench further stress API selection, argument generation, stateful tool use, and multi-turn user--tool interaction~\citep{li2023apibank,tang2023toolalpaca,qin2023toolllm,lu2024toolsandbox,yao2024taubench}. Together, these benchmarks motivate task-closure evaluation, where success depends on whether the environment reaches the intended final state rather than whether the model produces a plausible text answer alone.

Despite extensive studies on general LLMs, reasoning models, hallucination, autonomous agents, planning, memory, tool learning, and interactive benchmarks, there is still limited discussion of the boundary between the Agent Era and the OpenClaw Era~\citep{dong2024survey,huang2024understanding,yehudai2025survey,mamun2026anatomical}. Existing work provides many necessary ingredients, including external action loops, web and desktop environments, software repositories, stateful tools, execution feedback, memory mechanisms, reflection loops, and task-level verification~\citep{wei2025plangenllms,Rawles2024AndroidWorldAD,Bonatti2024WindowsAA,Kapoor2024OmniACTAD,Xu2024TheAgentCompanyBL,Merrill2026TerminalBenchBA,Boisvert2024WorkArenaTC}. However, these ingredients are often treated separately; in this paper, we re-examine them under a workspace-centered perspective~\citep{drouin2024workarena,chezelles2024browsergym,wang2024openhands,Boisvert2024WorkArenaTC,Xu2024TheAgentCompanyBL}. Specifically, we argue that files, terminals, browser sessions, logs, permissions, snapshots, and skill assets jointly define what an agent can perceive, modify, verify, and recover~\citep{Rawles2024AndroidWorldAD,Bonatti2024WindowsAA,Kapoor2024OmniACTAD,Merrill2026TerminalBenchBA}. This perspective clarifies why reliability, provenance, rollback, permissions, and governance become core architectural issues once agents operate over durable workspaces~\citep{zhao2026clawguard,ge2026governance}.

